\chapter{Introduzione}
\section{Contesto di ricerca}
I sistemi cyber-fisici combinano acquisizione, calcolo, comunicazione e attuazione. L'apprendimento multimodale pu\`o migliorare la consapevolezza situazionale, ma introduce anche questioni di robustezza, interpretabilit\`a, privacy e riproducibilit\`a \cite{goodfellow2016deep,nist2023airmf}.

\begin{researchquestion}[RQ1 -- Robustezza]
Come pu\`o un modello multimodale rimanere affidabile quando un flusso sensoriale \`e rumoroso, ritardato o non disponibile?
\end{researchquestion}

\begin{researchquestion}[RQ2 -- Fiducia]
Quali scelte progettuali e di valutazione rendono il comportamento del sistema pi\`u trasparente per operatori e auditor?
\end{researchquestion}

\section{Contributi}
\begin{keypoint}[Contributi]
\begin{enumerate}
  \item Un'architettura di fusione modulare per flussi sensoriali eterogenei.
  \item Una funzione obiettivo consapevole dell'affidabilit\`a e un protocollo per modalit\`a mancanti.
  \item Una suite di valutazione riproducibile per accuratezza, calibrazione, latenza e robustezza.
\end{enumerate}
\end{keypoint}

\chapter{Fondamenti e stato dell'arte}
\section{Apprendimento di rappresentazioni multimodali}
Organizzare lo stato dell'arte per domanda di ricerca, evitando un semplice catalogo di articoli. Confrontare ipotesi, dataset, metriche e limiti. Una tabella di sintesi pu\`o aiutare il lettore a comprendere la lacuna affrontata dalla tesi.

\begin{table}[htbp]
\centering
\caption{Confronto illustrativo tra approcci correlati.}
\begin{tabularx}{\textwidth}{Ycccc}
\toprule
Approccio & Fusione & Dati mancanti & Calibrazione & Adatto all'edge\\
\midrule
Fusione anticipata & Ingresso & Limitata & No & Parziale\\
Fusione tardiva & Decisione & Moderata & Opzionale & S\`i\\
Metodo proposto & Ibrida & Esplicita & S\`i & S\`i\\
\bottomrule
\end{tabularx}
\end{table}

\section{Intelligenza artificiale affidabile}
L'affidabilit\`a \`e multidimensionale. In questa tesi, le definizioni operative devono essere collegate a propriet\`a misurabili, come errore di calibrazione, robustezza alle perturbazioni, limiti di latenza e configurazioni sperimentali tracciabili.

\chapter{Metodologia}\label{ch:m-methodology}
\section{Architettura del sistema}
La Figura~\ref{fig:m-architecture} mostra una pipeline di ricerca leggera. Utilizza box LaTeX nativi invece di un motore grafico.

\begin{figure}[htbp]
\centering
\pipelinebox{Validazione della\\qualit\`a}\pipelinearrow
\pipelinebox{Encoder per\\modalit\`a}\pipelinearrow
\pipelinebox{Fusione basata\\sull'affidabilit\`a}\pipelinearrow
\pipelinebox{Predizione e\\spiegazione}
\caption{Architettura multimodale illustrativa.}
\label{fig:m-architecture}
\end{figure}

\section{Funzione obiettivo}
Una funzione obiettivo semplificata combina la perdita sul task, la regolarizzazione della calibrazione e la consistenza in presenza di dropout delle modalit\`a:
\begin{equation}
\mathcal{L}=\mathcal{L}_{\mathrm{task}}+\lambda_{c}\mathcal{L}_{\mathrm{cal}}+\lambda_{m}\mathcal{L}_{\mathrm{missing}}.
\end{equation}

\begin{keypoint}[Riproducibilit\`a]
Dichiarare versioni dei dataset, preprocessing, iperparametri, seed casuali, hardware, librerie software e criteri di arresto. Distinguere gli esperimenti esplorativi dalla valutazione confermativa.
\end{keypoint}

\chapter{Risultati e discussione}\label{ch:m-results}
\section{Risultati principali}
\begin{center}
\begin{minipage}{0.30\textwidth}\metric{94.2\%}{macro F1}{ISISBlue}\end{minipage}\hfill
\begin{minipage}{0.30\textwidth}\metric{0.031}{errore di calibrazione}{ISISLightBlue}\end{minipage}\hfill
\begin{minipage}{0.30\textwidth}\metric{27 ms}{latenza mediana}{ISISDeepBlue}\end{minipage}
\end{center}

\begin{figure}[htbp]
\centering
\begin{minipage}{0.84\textwidth}
{\sffamily\bfseries\color{ISISDeepBlue}Macro F1 con il 30\% di modalit\`a mancanti\par\vspace{3mm}}
\demobar{Baseline}{55mm}{69.0\%}{ISISMuted}
\demobar{Proposto}{68mm}{85.1\%}{ISISBlue}
\end{minipage}
\caption{Confronto di robustezza illustrativo senza PGFPlots.}
\label{fig:m-results}
\end{figure}

\begin{resultbox}[Risultato principale]
Il modello proposto conserva una quota maggiore delle prestazioni nominali al diminuire della disponibilit\`a delle modalit\`a. Sostituire questo testo con un'interpretazione statisticamente fondata degli esperimenti reali.
\end{resultbox}

\section{Minacce alla validit\`a}
Discutere validit\`a interna, di costrutto, esterna e delle conclusioni. Riportare risultati negativi e casi di fallimento quando migliorano la qualit\`a del resoconto scientifico.

\chapter{Conclusioni e prospettive}
Rispondere esplicitamente a ciascuna domanda di ricerca, riassumere i contributi originali e collegare i risultati al contesto scientifico pi\`u ampio. Separare le conclusioni supportate dalle evidenze dalle ipotesi per lavori futuri.

\section{Sviluppi futuri}
Possibili estensioni includono adattamento online, garanzie formali sull'incertezza, apprendimento rispettoso della privacy, sperimentazione sul campo e valutazione human-centred con gli operatori.
